The structural validity and predictive reliability of the Agent-Based Model (ABM) relied strictly on the empirical grounding of its initialization matrices. Before the Heuristic-Guided Deep Reinforcement Learning (HuDRL) agent could be trained to discover optimal policies, the underlying virtual environment had to accurately mirror the complex, heterogeneous reality of the Municipality of Bacolod. Real-world data gathered through legislative documents, household census records, and key informant interviews were quantitatively translated to strictly define the state space, operational constraints, and reward structures of the simulation environment. This meticulous parameterization ensured the AI was solving a genuine socioeconomic problem rather than an arbitrary mathematical puzzle.

\subsection{Empirical Parameterization of the State Space}

To guarantee that the DRL agent optimized within realistic and legally binding boundaries, the simulation's constraints were defined across three distinct environmental layers. This multi-layered parameterization mathematically prevented the AI from exploiting infinite resources or generating policies that would be legally impossible for the actual Local Government Unit (LGU) to execute.

\paragraph{Macro-Level Legislative and Financial Bounds:}
The extraction of municipal financial data from the Municipal Environment and Natural Resources Office (MENRO) and the codified Appropriation Ordinance No. 2024-01 established the absolute boundaries of the AI's action space. The virtual environment was strictly bounded by a \textpeso 1,500,000 annual Solid Waste Management (SWM) budget, which mathematically constrained the Mayor Agent's quarterly optimization limit ($B_q$) to an absolute maximum of \textpeso 375,000. 

Furthermore, operational constraints were dictated by regional economic realities. Incorporating the Region X Daily Minimum Wage of \textpeso 400 dictated that deploying a single municipal enforcer on a strict 30-day temporal contract consumed exactly \textpeso 12,000 per quarter. Computationally, this established a severe zero-sum economy within the simulation, creating a ``Personnel Trap.'' If the DRL agent attempted to maximize enforcement by consuming the entire budget on salaries, its action space for Information, Education, and Communication (IEC) campaigns or household incentives instantly collapsed to zero. This quantitative constraint accurately reflects the operational bottlenecks identified by \cite{Nishimura2022}, forcing the AI to learn strict algorithmic triage rather than relying on brute-force funding.

\paragraph{Meso-Level Spatial and Demographic Initialization:}
To capture the geographical and demographic reality of the municipality, the simulation grid instantiated 4,022 discrete \texttt{HouseholdAgent} objects based on exact municipal census data. The translation of qualitative interview data into quantitative arrays revealed profound spatial and economic heterogeneity across the seven constituent barangays. As detailed in Table \ref{tab:barangay_demographics}, this heterogeneity created a highly non-convex optimization landscape, meaning there was no single ``global'' policy that could succeed everywhere. 

\begin{table}[H]
    \centering
    \caption{Barangay Demographic, Financial, and Income Initialization Matrices}
    \label{tab:barangay_demographics}
    \small 
    \begin{tabular}{l r r c c c c}
        \toprule
        & & {Annual} & {Init.} & \multicolumn{3}{c}{{Income Profile (\%)}} \\
        \cmidrule(lr){5-7}
        {Barangay} & {Households} & {Budget (\textpeso)} & {Compliance} & {Low} & {Mid} & {High} \\
        \midrule
        Brgy Poblacion    & 1,534 & 200,000 & 2\%  & 70 & 25 & 5 \\
        Brgy Liangan East & 608   & 30,000  & 14\% & 40 & 40 & 20 \\
        Brgy Ezperanza    & 574   & 90,000  & 14\% & 20 & 50 & 30 \\
        Brgy Binuni       & 507   & 126,370 & 15\% & 50 & 30 & 20 \\
        Brgy Demologan    & 463   & 21,000  & 11\% & 80 & 15 & 5 \\
        Brgy Babalaya     & 171   & 15,000  & 15\% & 80 & 10 & 10 \\
        Brgy Mati         & 165   & 80,000  & 13\% & 90 & 5  & 5 \\
        \bottomrule
    \end{tabular}
    \vspace{1ex}
    
    \raggedright
    \footnotesize
    \textit{Note: Income profiles derived from Appendix B interviews. Low-income agents trigger a higher scalar multiplier ($\gamma$) for financial sensitivity within the ABM utility function. Initial Compliance reflects the calibrated state ($t=0$) to match empirical MENRO baseline data.}
\end{table}

The parameterization process identified three distinct socioeconomic typologies that fundamentally altered the marginal utility of money within the agents' decision-making algorithms:
\begin{itemize}
    \item \textbf{The Urban Resource Trap (Poblacion):} Despite possessing the highest absolute local budget (\textpeso 200,000), Poblacion suffered from the lowest per-capita spending power due to its massive population ($N=1,534$). Bound by its historical ``Police State'' profile, 60\% of its internal budget was locked into enforcement simply to maintain basic community order, mathematically capping its baseline autonomous compliance at an abysmal 2\%.
    \item \textbf{The Wealthy Anomaly (Binuni):} Operating with a robust \textpeso 126,370 budget for only 507 households, Binuni enjoyed the highest per-capita fiscal surplus. This demographic wealth generated an elastic response to policy, providing sufficient liquidity to effortlessly overcome high behavioral friction.
    \item \textbf{The Subsistence Trap (Babalaya \& Demologan):} Representing the economic floor of the municipality, 80\% of the households in these zones were classified as low-income. Operating on micro-budgets (\textpeso 15,000 to \textpeso 21,000), up to 90\% of their internal funds were consumed by bare-minimum enforcement honoraria. Computationally, this created a ``dead zone'' in the state space where autonomous IEC and incentive distributions were mathematically impossible without direct external municipal augmentation by the Mayor Agent.
\end{itemize}

\paragraph{Micro-Level Behavioral Matrices and the Intention-Action Gap:}
Beyond macro-economics, the simulation heavily parameterized the internal psychology of the citizens. The foundational behavioral parameters for the household agents were derived through a synthesis of empirical Knowledge, Attitude, and Practices (KAP) data \citep{Paigalan2025}. 

A critical finding from this data extraction was the presence of a severe ``Intention-Action Gap'' \citep{Geiger2021}. Consequently, agents were initialized with a relatively high mean theoretical Attitude ($A_0 \approx 0.66$) but drastically lower physical compliance rates. This mathematical dissonance forces the AI Mayor to recognize that citizens already \textit{know} they should segregate, but are blocked by physical friction; therefore, the AI must use its budget to alter the environment rather than just attempting to raise attitudes. Table \ref{tab:behavioral_parameters} outlines the specific Theory of Planned Behavior (TPB) weights that mathematically defined the unique behavioral friction of each spatial grid.

\begin{table}[H]
    \centering
    \caption{Calibrated Behavioral Weights and Base-Layer Allocation Strategies}
    \label{tab:behavioral_parameters}
    \resizebox{\textwidth}{!}{ 
    \begin{tabular}{l c c c c c c c c c}
        \toprule
        & \multicolumn{5}{c}{{Agent Behavior Parameters}} & & \multicolumn{3}{c}{{Base Budget Allocation (\%)}} \\
        \cmidrule(lr){2-6} \cmidrule(lr){8-10}
        {Barangay} & {Attitude ($w_a$)} & {Norms ($w_{sn}$)} & {Control ($PBC$)} & {Effort ($c_{\text{effort}}$)} & {Decay ($\delta$)} & & {Enf} & {Inc} & {IEC} \\
        \midrule
        Binuni       & 0.40 & 0.45 & 0.35 & 0.45 & 0.0010 & & 40 & 40 & 20 \\
        Ezperanza    & 0.20 & 0.35 & 0.25 & 0.65 & 0.0015 & & 50 & 30 & 20 \\
        Babalaya     & 0.35 & 0.50 & 0.35 & 0.40 & 0.0010 & & 90 & 5  & 5 \\
        Liangan East & 0.25 & 0.30 & 0.20 & 0.65 & 0.0035 & & 25 & 65 & 10 \\
        Poblacion    & 0.25 & 0.15 & 0.20 & 0.75 & 0.0030 & & 60 & 20 & 20 \\
        Demologan    & 0.35 & 0.35 & 0.30 & 0.55 & 0.0010 & & 85 & 10 & 5 \\
        Mati         & 0.30 & 0.40 & 0.25 & 0.60 & 0.0030 & & 30 & 50 & 20 \\
        \bottomrule
    \end{tabular}
    }
    \vspace{1ex}
    
    \raggedright
    \footnotesize
    \textit{Note: Agent behavior parameters correspond to TPB weights: $w_a$ (Attitude), $w_{sn}$ (Social Norms), $PBC$ (Perceived Behavior Control), $c_{\text{effort}}$ (Base Cost of Effort), and $\delta$ (Compliance Decay Rate). Budget allocation percentages reflect the hardcoded local strategy under the Status Quo scenario.}
\end{table}

The extreme variance in these parameters defined the unique vulnerabilities of each localized environment. For instance, Poblacion exhibited an ``Urban Isolation'' profile, characterized by an exceptionally low Social Norm weight ($w_{sn} = 0.20$). Computationally, this rendered spatial compliance feedback from Moore neighborhood queries mathematically ineffective at triggering social tipping points—meaning peer pressure could not save the city center. Conversely, Barangay Mati exhibited the highest motivation decay rate ($\delta = 0.10$), indicating that computational IEC investments were rapidly eroded as agrarian agents quickly reverted to traditional disposal practices.

\subsection{Baseline Calibration and Stochastic Seeding}
Before activating the DRL optimization policies, the simulation required a rigorous calibration phase to ensure it accurately reproduced the empirical ``Status Quo'' of Bacolod at $t=0$. The raw micro-level interview data—which falsely suggested a municipal compliance rate exceeding 50\%—had to be algorithmically corrected downward to match the objective, municipal-wide MENRO waste characterization audit ($\approx 12.5\%$). Table \ref{tab:calibration_data} details the analytical pivots applied to correct these inflated baseline vectors.

\begin{table}[htbp]
    \centering
    \caption{Self-Reported Compliance Rates vs.\ Calibrated Baseline Compliance Rates}
    \label{tab:calibration_data}
    \small
    \begin{tabular}{l c c p{6.5cm}}
        \toprule
        \textbf{Barangay Node} & \textbf{Self-Reported} & \textbf{Calibrated ($t=0$)} & \textbf{Primary Calibration Justification} \\
        \midrule
        Brgy Babalaya     & 100\% & 15.0\% & Adjusted for lack of persistent funding/monitoring. \\
        Brgy Binuni       & 95\%  & 15.0\% & Decoupled theoretical awareness from physical practice. \\
        Brgy Mati         & 70\%  & 13.0\% & Corrected for high habit decay rate in agrarian zones. \\
        Brgy Liangan East & 65\%  & 14.0\% & Smoothed transient project-based spikes (Eco-bricks). \\
        Brgy Demologan    & 60\%  & 11.0\% & Standardized to the municipal empirical baseline. \\
        Brgy Ezperanza    & 20\%  & 14.0\% & Aligned with the middle-class control group variance. \\
        Brgy Poblacion    & 2\%   & 2.0\%  & \textbf{Anchor Point} (Aligned perfectly with MENRO Audit). \\
        \bottomrule
    \end{tabular}
\end{table}

To reconcile the remaining micro-level bias with the macro-level ground truth, the initialized states for the non-anchor barangays were assigned stochastically using a bounded uniform distribution, $U(0.11, 0.15)$. The mathematical expected value ($\mu$) of this distribution tightly anchors to the 12.5\% MENRO baseline. By executing this randomization, headless simulation benchmarks successfully stabilized at a Global Average Compliance of $\approx 12.5\%$ at $t = 0$. This mathematically proved the structural validity of the ABM, ensuring the subsequent DRL agent was optimizing against a rigorously realistic, noisy scenario rather than a mathematically trivial or perfectly uniform one \citep{McCulloch2022}.

\subsection{The ``Act vs. Reality'' Gap and the Danger of Algorithmic Bias}
The systemic necessity of the aforementioned calibration phase exposes a profound sociological and computational challenge in municipal governance: the danger of Socially Desirable Responding (SDR) \citep{Geiger2021}. The data extracted during the initialization phase quantitatively captured the ``Act vs. Reality'' gap described by \cite{Yazawa2025}, highlighting how local officials frequently conflate theoretical policy awareness with actual, physical execution on the ground.

From a machine learning perspective, addressing this statistical divergence was paramount to preventing algorithmic bias. Had the simulation trusted the unverified self-reports (e.g., Binuni's 95\% or Babalaya's 100\%), the RL environment would have suffered from a ``garbage in, garbage out'' (GIGO) failure. The AI Mayor would have been trained on a trivialized state space, systematically learning to ignore or underfund these crucial sectors under the mathematically false assumption that they had already achieved permanent environmental sustainability. 

The analytical pivots required to construct Table \ref{tab:calibration_data} directly expose three critical administrative blind spots that cause real-world policies to fail:
\begin{enumerate}
    \item \textbf{The Intention-Action Trap (Binuni):} As empirical evidence suggests, high positive attitudes frequently coexist with poor R.A. 9003 compliance. Treating verbal support for environmental initiatives as equivalent to physical segregation leads human administrators to prematurely withdraw municipal funding before habits are truly locked in.
    \item \textbf{The Transient Project Trap (Liangan East):} Spikes in compliance driven by temporary initiatives (such as the barter of Eco-bricks for rice) do not represent deeply habituated behavioral norms. The calibration corrects the dangerous assumption that behavior will persist; in reality, when external funds vanish, the behavior instantly decays ($\delta$) back to baseline.
    \item \textbf{The Urban Resource Trap (Poblacion):} Strikingly, the only data point where qualitative reporting matched the objective quantitative audit was the massive commercial center (2\%). This powerfully validates the sociological theory that high population density and urban anonymity actively dissolve social pressure, creating an environment that defaults to non-compliance regardless of the existence of legislative frameworks \citep{Villanueva2021}.
\end{enumerate}

Ultimately, correcting this statistical divergence proved that a successful smart governance framework cannot rely on subjective, politically motivated municipal reporting. To successfully bridge the gap between legislative intent and physical action, the AI's optimization logic—and by extension, the LGU's real-world policy design—must be anchored strictly to objective, systemic ground-truths.